\chapter{Studi Literatur}

\section{Qubit}
\par
Dalam sistem komputasi klasik, bit merupakan satuan informasi paling dasar yang hanya dapat memiliki dua keadaan diskrit, yaitu 0 atau 1 \cite{nielsen2010quantum}. Setiap operasi logika dalam komputer klasik dibangun dari manipulasi terhadap bit-bit ini, yang secara fisik dapat direpresentasikan oleh dua kondisi berbeda, seperti arus listrik yang mengalir atau tidak mengalir, tegangan tinggi atau rendah, serta keadaan magnetik utara atau selatan \cite{tanenbaum2013structured}. Dengan demikian, bit berfungsi sebagai pembawa informasi biner yang menjadi fondasi bagi seluruh proses komputasi konvensional \cite{shannon1948mathematical}.
\par
Namun, seiring berkembangnya pemahaman terhadap prinsip-prinsip mekanika kuantum, muncul paradigma baru dalam penyimpanan dan pemrosesan informasi yang dikenal sebagai qubit (quantum bit) \cite{nielsen2010quantum}. Tidak seperti bit klasik yang hanya dapat berada dalam satu keadaan tertentu pada satu waktu, qubit mampu berada dalam superposisi antara keadaan $\ket{0}$ dan $\ket{1}$ secara bersamaan \cite{dirac1981principles}. Transisi dari bit ke qubit ini menandai pergeseran fundamental dari logika deterministik menuju logika probabilistik-kuantum, di mana informasi tidak lagi direpresentasikan secara biner murni, melainkan sebagai vektor dalam ruang Hilbert dua dimensi yang memiliki representasi matriks sebagai berikut \cite{sakurai2017modern}.

\begin{align*}
    \ket{0} = \begin{pmatrix} 1 \\ 0 \end{pmatrix}
\end{align*}
dan
\begin{align*}
    \ket{1} = \begin{pmatrix} 0 \\ 1 \end{pmatrix}
\end{align*}
\par
Kombinasi linier dari dua qubit atau basis vektor yang menyatakan keadaan informasi kuantum secara umum dituliskan menjadi
\begin{equation}
\ket{\psi} = \alpha \ket{0} + \beta \ket{1}
\end{equation}
dengan $\alpha$ dan $\beta$ adalah bilangan kompleks. Keadaan tersebut harus memenuhi syarat normalisasi atau $\langle \psi | \psi \rangle = 1$. Sehingga $\alpha$ dan $\beta$ harus memenuhi
\begin{align*}
|\alpha|^2 + |\beta|^2 = 1
\end{align*}

\par
Sistem yang terdiri atas lebih dari satu qubit maka keadaan sistemnya dapat dinyatakan dengan perkalian langsung dari masing-masing keadaannya. Misalkan informasi yang dimiliki yaitu bilangan 2 dan 6 yang apabila dinyatakan dalam bit klasik akan menjadi 10 dan 110. Informasi tersebut jika dinyatakan secara kuantum maka dapat dituliskan menjadi
\begin{align*}
10 = \ket{1} \otimes \ket{0}
\end{align*}
dan
\begin{align*}
110 = \ket{1} \otimes \ket{1} \otimes \ket{0}
\end{align*}
atau secara ringkas juga dapat dituliskan menjadi
\begin{align*}
\ket{1} \otimes \ket{0} = \ket{10}
\end{align*}
dan
\begin{align*}
\ket{1} \otimes \ket{1} \otimes \ket{0} = \ket{110}
\end{align*}
\par
Secara umum kombinasi linier untuk sistem dengan keadaan n qubit dapat dituliskan sebagai berikut
\begin{equation}
    \ket{\psi} = \ket{\psi_1}\otimes\ket{\psi_2}\otimes\ket{\psi_3}\otimes...\otimes\ket{\psi_n}
\end{equation}
Jika $\ket{\psi} = \alpha \ket{0} + \beta \ket{1}$, maka
\begin{equation}
    \ket{\psi} = (\alpha_1 \ket{0} + \beta_1 \ket{1})\otimes(\alpha_2 \ket{0} + \beta_2 \ket{1})\otimes(\alpha_3 \ket{0} + \beta_3 \ket{1})\otimes...\otimes(\alpha_n \ket{0} + \beta_n \ket{1})
\end{equation}
Persamaan umum untuk kubit rangkap dua dapat dituliskan sebagai berikut
\begin{align*}
    \ket{\psi} = \ket{\psi_1}\otimes\ket{\psi_2}
\end{align*}
sehingga
\begin{align*}
    \ket{\psi} = \alpha_1\alpha_2\ket{00} + \alpha_1\beta_2\ket{01} + \alpha_2\beta_1\ket{10} + \beta_1\beta_2\ket{11}
\end{align*}
atau bisa ditulis dengan
\begin{equation}
    \ket{\psi} = c_1\ket{00} + c_2\ket{01} + c_3\ket{10} + c_4\ket{11}
\end{equation}
dengan $\alpha_1\alpha_2=c_1$, $\alpha_1\beta_2=c_2$, $\beta_1\alpha_2=c_3$, dan $\beta_1\beta_2=c_4$, di mana d merupakan bilangan kompleks sehingga sistem kuantum dengan qubit rangkap dua harus ternormalisasi dan harus memenuhi
\begin{equation}
    |d_1|^2+|d_2|^2+|d_3|^2+|d_4|^2=1
\end{equation}
\par
Persamaan umum untuk keadaan kubit rangkap dua tersebut jika disederhanakan dalam bentuk notasi sigma adalah sebagai berikut
\begin{equation}
    \ket{\psi}=\sum_{i,j=0}^{1} d_{ij} \ket{ij}
\end{equation}
\par
Qubit juga memiliki realitas fisis, beberapa diantaranya adalah pada kasus polarisasi \cite{nielsen2010quantum}. Pada polarisasi, $\ket{0}$ merupakan representasi dari polarisasi vertikal foton, sedangkan $\ket{1}$ untuk polarisasi horizontalnya \cite{kok2007linear}. Kubit juga dapat digunakan menjadi representasi dari spin elektron, misalnya $\ket{0}$ adalah representasi dari spin-up dan $\ket{1}$ sebagai spin-down \cite{sakurai2017modern}.

\section{Matriks Densitas}
\par
Ungkapan matematika lain untuk informasi sistem fisis direpresentasikan oleh matriks densitas. Matriks densitas dari campuran n keadaan murni dengan probabilitas $p_i$ dapat dituliskan sebagai 
\begin{equation}\label{rho}
	\rho=\sum_{i=1}^{n}p_i\ket{\psi_i}\bra{\psi_i},
\end{equation}
dengan
\begin{equation}
	\sum_{i=1}^{n}p_i=1,
\end{equation}
di mana matriks $\rho$ harus bersifat berikut:
\begin{enumerate}
	\item \text{Hermitian}
	\begin{equation}
		\rho=\rho^\dag,
	\end{equation}
	\item \text{Trace dari matriks densitas bernilai 1}
	\begin{equation}
		\text{tr}(\rho)=1,
	\end{equation}
	\item \text{Semua nilai eigennya bernilai non-negatif \cite{artawan2022invariants}}.
	\begin{equation}
		\bra{x}\rho\ket{x}\geq 0.
	\end{equation}
	\par
\end{enumerate}
\par
Untuk kasus di mana $\rho$ merepresentasikan 1 keadaan murni (pure state), nilai $i$ hanya 1, sehingga persamaan (\ref{rho}) menjadi
\begin{equation}\label{aaaaa}
	\rho=\ket{\psi}\bra{\psi}.
\end{equation}
Berlaku $Tr(\rho^2)=1$ dan $\rho^2=\rho$ \cite{nakahara2008quantum}.
\par
Menggunakan representasi pada Bola Bloch, di mana  $\ket{\psi} = \alpha \ket{0} + \beta \ket{1}$ diganti menjadi $\ket{\psi} =\cos\theta\ket{0}+\sin\theta e^{i\phi}\ket{1}$ serta persamaan (\ref{aaaaa}), matriks densitas dari keadaan satu kubit $\ket{\psi}$ dapat diturunkan sebagai berikut
	\begin{align}\label{rho1kubit}
		\rho&=(\cos\theta\ket{0}+\sin\theta e^{i\phi}\ket{1})(\cos\theta\bra{0}+\sin\theta e^{-i\phi}\bra{1})\nonumber\\
		&=\begin{pmatrix}
			\cos^2\theta & \sin\theta\cos\theta e^{-i\phi} \\
			\sin\theta\cos\theta e^{i\phi} & \sin^2\theta
		\end{pmatrix}\nonumber\\
		&=\begin{pmatrix}
			\frac{1}{2}(\cos2\theta+1)& \sin\theta\cos\theta(\cos\phi-i\sin\phi)\\
			\sin\theta\cos\theta(\cos\phi+i\sin\phi)& \frac{1}{2}(1-\cos2\theta)
		\end{pmatrix}\nonumber\\
		&=\frac{1}{2}\left\{\begin{pmatrix}
			1 & 0 \\
			0 & 1
		\end{pmatrix}+\sin2\theta\cos\phi\begin{pmatrix}
			0 & 1 \\
			1 & 0
		\end{pmatrix}+\sin2\theta\sin\phi\begin{pmatrix}
			0 & -i\\
			i & 0
		\end{pmatrix}+ \cos2\theta\begin{pmatrix}
			1 & 0 \\
			0 & -1
		\end{pmatrix}\right\}\nonumber\\			&=\frac{1}{2}\left\{I+n_1\sigma_1+n_2\sigma_2+n_3\sigma_3\right\}\nonumber\\
		&=\frac{1}{2}(I+\hat{n}\cdot\vec{\sigma}),
	\end{align}
di mana $\hat{n}=\sin2\theta\cos\phi\hat{i}+\sin2\theta\sin\phi\hat{j}+\cos2\theta\hat{k}$ dan $\vec{\sigma}=\sigma_1\hat{i}+\sigma_2\hat{j}+\sigma_3\hat{k}$ dengan $\sigma_1$, $\sigma_2$ dan $\sigma_3$ merupakan matriks Pauli.

\section{Matriks Densitas Tereduksi}

Jika terdapat suatu matriks persegi atau operator $T$, maka trace dari matriks $T$ tersebut adalah jumlahan elemen diagonalnya, atau bisa dituliskan sebagai

\begin{equation}
\mathrm{tr}(T) \equiv \sum_i \langle i | T | i \rangle , 
\end{equation}

dimana $|i\rangle$ adalah suatu basis orthonormal. Misalkan ditinjau matriks densitas sistem qubit ganda $\rho_{AB} = \rho_A \otimes \rho_B$, dan kita hanya ingin mengetahui salah satu sub-sistem saja, misalkan sub-sistem $A$. Maka kita bisa menghilangkan atau trace sub-sistem $B$ saja, sehingga akan menyisakan matriks densitas tereduksi untuk sub-sistem $A$

\begin{equation}
\rho_A = \mathrm{tr}_B(\rho_{AB})
\end{equation}

\begin{equation}\\
= \sum_{m=0}^{1} \rho_A \otimes \langle m| \rho_B |m\rangle
\end{equation}

\begin{equation}
\rho_A =
\sum_{m=0}^{1}
(I_A \otimes \langle m|_B)\, \rho_{AB}\, (I_A \otimes |m\rangle_B) .
\end{equation}

Sama halnya, jika kita ingin meninjau sub-sistem $B$, maka kita bisa menghilangkan sub-sistem $A$ menjadi (Vathsan, 2015)

\begin{equation}
\rho_B =
\sum_{m=0}^{1}
(\langle m|_A \otimes I_B)\, \rho_{AB}\, (|m\rangle_A \otimes I_B) .
\end{equation}

\subsection{Matriks Densitas Tereduksi Dua Qubit}

Pada keadaan dua qubit, matriks densitas tereduksi yang mungkin hanya $\rho_A$ dan $\rho_B$, karena hanya terdiri atas sub-sistem $A$ dan sub-sistem $B$ saja. Untuk matriks densitas tereduksi sub-sistem $A$ secara umum bisa dituliskan sebagai

\begin{align}
\rho_A &= \mathrm{tr}_B(\rho_{AB}) \nonumber \\
&= \sum_{m=0}^{1} (I \otimes \langle m|)\, \rho_{AB}\, (I \otimes |m\rangle) \nonumber \\
&= \sum_{m=0}^{1} \sum_{i_1,j_1=0}^{1} \sum_{i_2,j_2=0}^{1}
\alpha_{i_1 j_1}\alpha^{*}_{i_2 j_2}
(I \otimes \langle m|)\, |i_1 j_1\rangle \langle i_2 j_2|\, (I \otimes |m\rangle) \nonumber \\
&= \sum_{m=0}^{1} \sum_{i_1,j_1=0}^{1} \sum_{i_2,j_2=0}^{1}
\alpha_{i_1 j_1}\alpha^{*}_{i_2 j_2}
I |i_1\rangle \langle m|j_1\rangle \langle i_2| I \langle j_2|m\rangle \nonumber \\
&= \sum_{m=0}^{1} \sum_{i_1,j_1=0}^{1} \sum_{i_2,j_2=0}^{1}
\alpha_{i_1 j_1}\alpha^{*}_{i_2 j_2}
\delta_{m j_1}\delta_{j_2 m} |i_1\rangle \langle i_2| \nonumber \\
\rho_A &= \sum_{m=0}^{1} \sum_{i_1,i_2=0}^{1}
\alpha_{i_1 m}\alpha^{*}_{i_2 m}
|i_1\rangle \langle i_2| .
\end{align}

sedangkan matriks tereduksi sub-sistem $B$ secara umum bisa dituliskan sebagai

\begin{align}
\rho_B &= \mathrm{tr}_A(\rho_{AB}) \nonumber \\
&= \sum_{m=0}^{1} (\langle m| \otimes I)\, \rho_{AB}\, (|m\rangle \otimes I) \nonumber \\
&= \sum_{m=0}^{1} \sum_{i_1,j_1=0}^{1} \sum_{i_2,j_2=0}^{1}
\alpha_{i_1 j_1}\alpha^{*}_{i_2 j_2}
\langle m|i_1\rangle I |j_1\rangle \langle i_2|m\rangle \langle j_2| I \nonumber \\
\rho_B &= \sum_{m=0}^{1} \sum_{j_1,j_2=0}^{1}
\alpha_{m j_1}\alpha^{*}_{m j_2}
|j_1\rangle \langle j_2| .
\end{align}

\subsection{Matriks Densitas Tereduksi Tiga Qubit}

Matriks densitas tereduksi pada tiga qubit terdiri dari sub-sistem qubit tunggal dan sub-sistem qubit ganda. Untuk sub-sistem qubit tunggal terdiri dari $\rho_A$, $\rho_B$, dan $\rho_C$, sedangkan untuk sub-sistem qubit ganda terdiri dari $\rho_{AB}$, $\rho_{AC}$, dan $\rho_{BC}$. Untuk sub-sistem qubit tunggal, matriks densitas tereduksi sub-sistem $A$ secara umum dapat dituliskan sebagai (Yuwana, 2018)

\begin{align}
\rho_A &= \mathrm{tr}_{BC}(\rho_{ABC}) \nonumber \\
&= \sum_{m,n=0}^{1} (I \otimes \langle m| \otimes \langle n|)\, \rho_{ABC}\, (I \otimes |m\rangle \otimes |n\rangle) \nonumber \\
&= \sum_{m,n=0}^{1} \sum_{i_1,j_1,k_1=0}^{1} \sum_{i_2,j_2,k_2=0}^{1}
\alpha_{i_1 j_1 k_1}\alpha^{*}_{i_2 j_2 k_2}
I |i_1\rangle \langle mn|j_1 k_1\rangle \langle i_2| I \langle j_2 k_2|mn\rangle \nonumber \\
\rho_A &= \sum_{m,n=0}^{1} \sum_{i_1,i_2=0}^{1}
\alpha_{i_1 m n}\alpha^{*}_{i_2 m n}
|i_1\rangle \langle i_2| ,
\end{align}

matriks densitas tereduksi sub-sistem $B$ secara umum dapat dituliskan sebagai

\begin{align}
\rho_B &= \mathrm{tr}_{AC}(\rho_{ABC}) \nonumber \\
&= \sum_{m,n=0}^{1} (\langle m| \otimes I \otimes \langle n|)\, \rho_{ABC}\, (|m\rangle \otimes I \otimes |n\rangle) \nonumber \\
&= \sum_{m,n=0}^{1} \sum_{i_1,j_1,k_1=0}^{1} \sum_{i_2,j_2,k_2=0}^{1}
\alpha_{i_1 j_1 k_1}\alpha^{*}_{i_2 j_2 k_2}
\langle m|i_1\rangle I |j_1\rangle \langle n|k_1\rangle
\langle m|i_2\rangle I |j_2\rangle \langle n|k_2\rangle \nonumber \\
\rho_B &= \sum_{m,n=0}^{1} \sum_{j_1,j_2=0}^{1}
\alpha_{m j_1 n}\alpha^{*}_{m j_2 n}
|j_1\rangle \langle j_2| ,
\end{align}

dan matriks densitas tereduksi sub-sistem $C$ secara umum dapat dituliskan sebagai

\begin{align}
\rho_C &= \mathrm{tr}_{AB}(\rho_{ABC}) \nonumber \\
&= \sum_{m,n=0}^{1} (\langle m| \otimes \langle n| \otimes I)\, \rho_{ABC}\, (|m\rangle \otimes |n\rangle \otimes I) \nonumber \\
&= \sum_{m,n=0}^{1} \sum_{i_1,j_1,k_1=0}^{1} \sum_{i_2,j_2,k_2=0}^{1}
\alpha_{i_1 j_1 k_1}\alpha^{*}_{i_2 j_2 k_2}
\langle mn|i_1 j_1\rangle I |k_1\rangle
\langle i_2 j_2|mn\rangle \langle k_2| I \nonumber \\
\rho_C &= \sum_{m,n=0}^{1} \sum_{k_1,k_2=0}^{1}
\alpha_{m n k_1}\alpha^{*}_{m n k_2}
|k_1\rangle \langle k_2| .
\end{align}

Untuk sub-sistem qubit ganda, matriks densitas tereduksi sub-sistem $AB$ secara umum dapat dituliskan sebagai

\begin{align}
\rho_{AB} &= \mathrm{tr}_C(\rho_{ABC}) \nonumber \\
&= \sum_{m=0}^{1} (I \otimes I \otimes \langle m|)\, \rho_{ABC}\, (I \otimes I \otimes |m\rangle) \nonumber \\
&= \sum_{m=0}^{1} \sum_{i_1,j_1,k_1=0}^{1} \sum_{i_2,j_2,k_2=0}^{1}
\alpha_{i_1 j_1 k_1}\alpha^{*}_{i_2 j_2 k_2}
I |i_1\rangle I |j_1\rangle \langle m|k_1\rangle
\langle i_2| I \langle j_2| I \langle k_2|m\rangle \nonumber \\
\rho_{AB} &= \sum_{m=0}^{1} \sum_{i_1,j_1=0}^{1} \sum_{i_2,j_2=0}^{1}
\alpha_{i_1 j_1 m}\alpha^{*}_{i_2 j_2 m}
|i_1 j_1\rangle \langle i_2 j_2| ,
\end{align}

matriks densitas tereduksi sub-sistem $AC$ secara umum dapat dituliskan sebagai

\begin{align}
\rho_{AC} &= \mathrm{tr}_B(\rho_{ABC}) \nonumber \\
&= \sum_{m=0}^{1} (I \otimes \langle m| \otimes I)\, \rho_{ABC}\, (I \otimes |m\rangle \otimes I) \nonumber \\
&= \sum_{m=0}^{1} \sum_{i_1,j_1,k_1=0}^{1} \sum_{i_2,j_2,k_2=0}^{1}
\alpha_{i_1 j_1 k_1}\alpha^{*}_{i_2 j_2 k_2}
I |i_1\rangle \langle m|j_1\rangle I |k_1\rangle
\langle i_2| I \langle j_2|m\rangle \langle k_2| I \nonumber \\
\rho_{AC} &= \sum_{m=0}^{1} \sum_{i_1,k_1=0}^{1} \sum_{i_2,k_2=0}^{1}
\alpha_{i_1 m k_1}\alpha^{*}_{i_2 m k_2}
|i_1 k_1\rangle \langle i_2 k_2| ,
\end{align}

dan matriks densitas tereduksi sub-sistem $BC$ secara umum dapat dituliskan sebagai

\begin{align}
\rho_{BC} &= \mathrm{tr}_A(\rho_{ABC}) \nonumber \\
&= \sum_{m=0}^{1} (\langle m| \otimes I \otimes I)\, \rho_{ABC}\, (|m\rangle \otimes I \otimes I) \nonumber \\
&= \sum_{m=0}^{1} \sum_{i_1,j_1,k_1=0}^{1} \sum_{i_2,j_2,k_2=0}^{1}
\alpha_{i_1 j_1 k_1}\alpha^{*}_{i_2 j_2 k_2}
\langle m|i_1\rangle I |j_1\rangle I |k_1\rangle
\langle i_2|m\rangle \langle j_2| I \langle k_2| I \nonumber \\
\rho_{BC} &= \sum_{m=0}^{1} \sum_{j_1,k_1=0}^{1} \sum_{j_2,k_2=0}^{1}
\alpha_{m j_1 k_1}\alpha^{*}_{m j_2 k_2}
|j_1 k_1\rangle \langle j_2 k_2| .
\end{align}


\begin{equation}
|\psi\rangle_{\mathrm{GHZ}} = \frac{1}{\sqrt{2}} (|000\rangle + |111\rangle) .
\end{equation}

\section{Prinsip Dasar Keadaan Terbelit}
\par
Suatu keadaan murni $\ket{\psi_{12}}$ dikatakan terbelit jika dan hanya jika tidak memenuhi
\begin{equation}
	\ket{\psi_{12}}=\ket{\psi_1}\otimes\ket{\psi_2}
\end{equation}
atau dikatakan sebagai keadaan tidak terpisah (\textit{non-separable}). Jika memenuhi maka dikatakan keadaan terpisah (\textit{separable}) \cite{esteve2004quantum}.
\par
Untuk memahami pernyataan di atas, dimisalkan $\ket{\psi_1}$ dan $\ket{\psi_2}$ dalam keadaan dua kubit sebagai berikut
\begin{equation}
	\begin{aligned}
		\ket{\psi_1}=x_0\ket{0}+x_1\ket{1}\\
		\ket{\psi_2}=y_0\ket{0}+y_1\ket{1}
	\end{aligned}
\end{equation}
dan suatu keadaan dua kubit yang belum diketahui keterbelitannya
\begin{equation}
	\begin{aligned}
		\ket{\psi_{12}}&=z_{00}\ket{00}+z_{01}\ket{01}+z_{10}\ket{10}+z_{11}\ket{11}
	\end{aligned}
\end{equation}
di mana
\begin{align}
    z_{00}=x_0y_0
    \label{eq:terbelit1}\\
    z_{01}=x_0y_1
    \label{eq:terbelit2}\\
    z_{10}=x_1y_0
    \label{eq:terbelit3}\\
    z_{11}=x_1y_1
    \label{eq:terbelit4}
\end{align}
Jika koefisien $z_{00}$, $z_{01}$, $z_{10}$, dan $z_{11}$ berturut-turut bisa diuraikan dalam perkalian $x_0y_0$, $x_0y_1$, $x_1y_0$, dan $x_1y_1$, maka keadaan dikatakan terpisah. Jika tidak bisa diuraikan dalam perkalian tersebut, maka keadaan dikatan terbelit \cite{taufiqi2024asymmetric}.
\par
Sebagai contoh diberikan keadaan sebagai berikut
\begin{equation}
	\ket{\psi}=\frac{1}{2}\left(\ket{00}-\ket{01}+\ket{10}-\ket{11}\right).
\end{equation}
Maka koefisien basisnya 
\begin{equation}
	z_{00}=-z_{01}=z_{10}=-z_{11}=\frac{1}{2}.
\end{equation}
Karena belum diketahui terbelit atau tidak, koefisien basis $z_{00}$, $z_{01}$, $z_{10}$, dan $z_{11}$ berturut-turut kita anggap terlebih dahulu sebagai perkalian $x_0y_0$, $x_0y_1$, $x_1y_0$, dan $x_1y_1$. Selanjutnya untuk mengetahui keterbelitannya, dilakukan pembagian antara koefisien basis sebagai berikut
\begin{equation}
	\frac{z_{00}}{z_{11}}=\frac{x_0y_0}{x_1y_0}=\frac{x_0}{x_1}=1.
\end{equation}
Dari perhitungan di atas didapatkan $x_0=x_1=\frac{1}{\sqrt{2}}$ yang memenuhi normalisasi $|x_0|^2+|x_1|^2=1$. Lalu untuk pembagian yang lain yaitu
\begin{equation}
	\frac{z_{00}}{z_{01}}=\frac{x_0y_0}{x_0y_1}=\frac{y_0}{y_1}=-1.
\end{equation}
Dengan diketahuinya $x_0$ dan $z_{00}$ maka $y_1=-y_0=-\frac{1}{\sqrt{2}}$ dan memenuhi $|y_0|^2+|y_1|^2=1$. Keadaan ini disebut keadaan terpisah sebab diketahui nilai $x_0$, $x_1$, $y_0$ dan $y_1$ yang artinya $z_{00}$, $z_{01}$, $z_{10}$, dan $z_{11}$ dapat diuraikan dalam perkalian koefisien dari 2 keadaan penyusunnya. Oleh karenanya keadaan 2 kubit dapat dituliskan sebagai
\begin{equation}
	\begin{aligned}
		\ket{\psi}&=\ket{\psi_1}\otimes\ket{\psi_2}\\
		&=\frac{1}{\sqrt{2}}\left(\ket{0}+\ket{1}\right)\frac{1}{\sqrt{2}}\left(\ket{0}-\ket{1}\right)
	\end{aligned}
\end{equation}
\par
Untuk keadaan dua kubit yang terbelit, maka nilai $x_0$, $x_1$, $y_0$ dan $y_1$ tidak ditemukan dengan cara di atas. Artinya koefisien basis $z_{00}$, $z_{01}$, $z_{10}$, dan $z_{11}$ berturut-turut tidak bisa diuraikan dalam perkalian $x_0y_0$, $x_0y_1$, $x_1y_0$, dan $x_1y_1$. Contoh dari keadaan 2 kubit yang terbelit yaitu 4 keadaan Bell \cite{zeilinger1998entanglement}
\begin{equation}
	\begin{aligned}
		\ket{\psi_1}&=\frac{1}{\sqrt{2}}\left(\ket{00}+\ket{11}\right)\\
		\ket{\psi_2}&=\frac{1}{\sqrt{2}}\left(\ket{00}-\ket{11}\right)\\
		\ket{\psi_3}&=\frac{1}{\sqrt{2}}\left(\ket{10}+\ket{01}\right)\\
		\ket{\psi_4}&=\frac{1}{\sqrt{2}}\left(\ket{01}-\ket{10}\right)
	\end{aligned}
\end{equation}

\section{Pola Keterbelitan dan Klasifikasi Keterbelitan}

\subsection{Koefisien Basis}

\subsubsection{Keadaan Terpisah}

Ditinjau contoh keadaan dua qubit $|\psi_1\rangle$ sebagai berikut

\begin{equation}
|\psi_1\rangle =
\frac{1}{2}
\left(
|00\rangle - |01\rangle + |10\rangle - |11\rangle
\right).
\end{equation}

Keadaan $|\psi_1\rangle$ dapat dikatakan sebagai keadaan yang terpisah jika persamaan mempunyai solusi yang pasti untuk nilai amplitudo probabilitas atau koefisien vektor
basis nya, yaitu

\begin{equation}\label{1}
\frac{1}{2} = a_0 b_0
\end{equation}

\begin{equation}\label{2}
-\frac{1}{2} = a_0 b_1
\end{equation}

\begin{equation}\label{3}
\frac{1}{2} = a_1 b_0
\end{equation}

\begin{equation}\label{4}
-\frac{1}{2} = a_1 b_1.
\end{equation}

Untuk mencari nilai $a_0$ dan $a_1$, kita bisa membagi persamaan (\ref{1}) dengan persamaan (\ref{4})

\begin{equation}
1 = \frac{a_0}{a_1}
\end{equation}

\begin{equation}
a_0 = a_1,
\end{equation}

dan karena syarat normalisasi, maka

\begin{equation}
a_0 = \frac{1}{\sqrt{2}} ; \quad
a_1 = \frac{1}{\sqrt{2}}.
\end{equation}

Cara lain mencari nilai $a_0$ dan $a_1$ adalah dengan membagi persamaan (\ref{2}) dengan persamaan (\ref{4})

\begin{equation}
1 = \frac{a_0}{a_1}
\end{equation}

\begin{equation}
a_0 = a_1
\end{equation}

dan karena syarat normalisasi, maka

\begin{equation}
a_0 = \frac{1}{\sqrt{2}} ; \quad
a_1 = \frac{1}{\sqrt{2}}.
\end{equation}

Untuk mencari nilai $b_0$ dan $b_1$, kita bisa membagi persamaan (\ref{1}) dengan persamaan (\ref{2})

\begin{equation}
-1 = \frac{b_0}{b_1}
\end{equation}

\begin{equation}
b_0 = - b_1
\end{equation}

\begin{equation}
b_0 = \frac{1}{\sqrt{2}} ; \quad
b_1 = -\frac{1}{\sqrt{2}}.
\end{equation}

Cara lain mencari nilai $b_0$ dan $b_1$ adalah dengan membagi persamaan (\ref{3}) dengan persamaan (\ref{4})

\begin{equation}
-1 = \frac{b_0}{b_1}
\end{equation}

\begin{equation}
b_0 = - b_1
\end{equation}

\begin{equation}
b_0 = \frac{1}{\sqrt{2}} ; \quad
b_1 = -\frac{1}{\sqrt{2}}.
\end{equation}

Karena hasil pada persamaan $a_0$ dan $a_1$ serta $b_0$ dan $b_1$konsisten menggunakan cara apapun, maka keadaan $|\psi_1\rangle$ memiliki solusi yang pasti untuk nilai koefisien $a_0, a_1, b_0,$ dan $b_1$,
sehingga keadaan $|\psi_1\rangle$ adalah keadaan yang terpisah.

\begin{equation}
|\psi_1\rangle =
\frac{1}{\sqrt{2}}
\left(
|0\rangle + |1\rangle
\right)
\otimes
\frac{1}{\sqrt{2}}
\left(
|0\rangle - |1\rangle
\right).
\end{equation}

\subsubsection{Keadaan Terbelit}

Ditinjau contoh keadaan yang lain $|\psi_2\rangle$ sebagai berikut

\begin{equation}
|\psi_2\rangle =
\frac{1}{2}
\left(
|00\rangle + |01\rangle + |10\rangle - |11\rangle
\right).
\end{equation}

Nilai koefisiennya adalah

\begin{equation}
\frac{1}{2} = a_0 b_0
\end{equation}

\begin{equation}
\frac{1}{2} = a_0 b_1
\end{equation}

\begin{equation}
\frac{1}{2} = a_1 b_0
\end{equation}

\begin{equation}
-\frac{1}{2} = a_1 b_1.
\end{equation}

Dengan cara yang sama dengan contoh keadaan $|\psi_1\rangle$, diperoleh hasil yang tidak
konsisten sehingga tidak ada solusi yang pasti untuk nilai koefisien $a_0, a_1, b_0,$ dan
$b_1$. Oleh karena itu, keadaan $|\psi_2\rangle$ adalah keadaan terbelit.

\subsubsection{Keadaan Tiga Qubit}

Ditinjau keadaan tiga qubit secara umum $|\psi_{ABC}\rangle$ sebagai berikut

\begin{equation}
|\psi_{ABC}\rangle
=
\sum_{i,j,k=0}^{1}
\alpha_{ijk} |ijk|
\end{equation}
\begin{equation}
=
\alpha_{000} |000\rangle
+ \alpha_{001} |001\rangle
+ \alpha_{010} |010\rangle
+ \alpha_{011} |011\rangle
\end{equation}
\begin{equation}
\quad
+ \alpha_{100} |100\rangle
+ \alpha_{101} |101\rangle
+ \alpha_{110} |110\rangle
+ \alpha_{111} |111\rangle ,
\end{equation}

dengan syarat normalisasi

\begin{equation}
\sum_{i,j,k=0}^{1}
|\alpha_{ijk}|^2
=
|\alpha_{000}|^2
+ |\alpha_{001}|^2
+ \cdots
+ |\alpha_{110}|^2
+ |\alpha_{111}|^2
= 1.
\end{equation}

Seperti halnya keadaan dua qubit secara umum, keadaan tiga qubit ini bisa tersusun atas
tiga keadaan sub-sistem qubit tunggal $|\phi_A\rangle$, $|\phi_B\rangle$, dan $|\phi_C\rangle$
sebagai berikut

\begin{equation}
|\phi_A\rangle = a_0 |0\rangle + a_1 |1\rangle
\end{equation}
\begin{equation}
|\phi_B\rangle = b_0 |0\rangle + b_1 |1\rangle
\end{equation}
\begin{equation}
|\phi_C\rangle = c_0 |0\rangle + c_1 |1\rangle ,
\end{equation}

dengan keadaan $|\phi_A\rangle \in \mathcal{H}_A = \mathbb{C}^2$, keadaan
$|\phi_B\rangle \in \mathcal{H}_B = \mathbb{C}^2$, dan keadaan
$|\phi_C\rangle \in \mathcal{H}_C = \mathbb{C}^2$, sehingga keadaan
$|\psi_{ABC}\rangle$ berada dalam ruang Hilbert
$\mathcal{H}_A \otimes \mathcal{H}_B \otimes \mathcal{H}_C
= \mathbb{C}^2 \otimes \mathbb{C}^2 \otimes \mathbb{C}^2$.

\subsubsection{Keadaan Terpisah Sepenuhnya}

Keadaan umum dapat dinamakan keadaan yang terpisah sepenuhnya
(\textit{completely separable state}) jika keadaan tersebut dapat dinyatakan dalam
perkalian tensor antar tiga keadaan qubit tunggal nya

\begin{equation}
|\psi\rangle_{ABC}
=
|\phi\rangle_A \otimes |\phi\rangle_B \otimes |\phi\rangle_C
\end{equation}
\begin{equation}
=
a_0 b_0 c_0 |000\rangle
+ a_0 b_0 c_1 |001\rangle
+ a_0 b_1 c_0 |010\rangle
+ a_0 b_1 c_1 |011\rangle
\end{equation}
\begin{equation}
\quad
+ a_1 b_0 c_0 |100\rangle
+ a_1 b_0 c_1 |101\rangle
+ a_1 b_1 c_0 |110\rangle
+ a_1 b_1 c_1 |111\rangle ,
\end{equation}

dengan syarat normalisasi

\begin{equation}
\sum_{i=0}^{1} |a_i|^2
=
\sum_{j=0}^{1} |b_j|^2
=
\sum_{k=0}^{1} |c_k|^2
= 1 .
\end{equation}

Maka dari itu, kita bisa dapatkan

\begin{align}
\alpha_{000} &= a_0 b_0 c_0 \nonumber \\
\alpha_{001} &= a_0 b_0 c_1 \nonumber \\
\alpha_{010} &= a_0 b_1 c_0 \nonumber \\
\alpha_{011} &= a_0 b_1 c_1 \nonumber \\
\alpha_{100} &= a_1 b_0 c_0 \nonumber \\
\alpha_{101} &= a_1 b_0 c_1 \nonumber \\
\alpha_{110} &= a_1 b_1 c_0 \nonumber \\
\alpha_{111} &= a_1 b_1 c_1 .
\end{align}

Suatu keadaan tiga qubit dinamakan terpisah sepenuhnya jika kita bisa mencari solusi
nilai $a_0$, $b_0$, $c_0$, $a_1$, $b_1$ dan $c_1$. Contoh keadaan
yang terpisah sepenuhnya adalah sebagai berikut

\begin{equation}
|\chi_1\rangle =
\frac{1}{2}
\left(
|000\rangle + |001\rangle + |010\rangle + |011\rangle
\right) .
\end{equation}

Seperti halnya pada keadaan dua qubit, untuk mencari solusi,
kita bisa melakukan perhitungan sebagai berikut

\begin{equation}
\frac{a_1 b_0 c_0}{a_0 b_0 c_0}
=
\frac{0}{1/2}
\end{equation}

\begin{equation}
\frac{a_0 b_1 c_0}{a_0 b_0 c_0}
=
\frac{1/2}{1/2}
\end{equation}

\begin{equation}
\frac{a_0 b_0 c_1}{a_0 b_0 c_0}
=
\frac{1/2}{1/2}
\end{equation}

\[
\therefore \ a_1 = 0 ; \ a_0 = 1
\qquad
\therefore \ b_0 = b_1 = \frac{1}{\sqrt{2}}
\]

\[
\therefore \ c_0 = c_1 = \frac{1}{\sqrt{2}}
\]

\begin{equation}
\frac{a_1 b_0 c_1}{a_0 b_0 c_1}
=
\frac{0}{1/2}
\end{equation}

\begin{equation}
\frac{a_0 b_1 c_1}{a_0 b_0 c_1}
=
\frac{1/2}{1/2}
\end{equation}

\begin{equation}
\frac{a_0 b_1 c_1}{a_0 b_1 c_0}
=
\frac{1/2}{1/2}
\end{equation}

\[
\therefore \ a_1 = 0 ; \ a_0 = 1
\qquad
\therefore \ b_0 = b_1 = \frac{1}{\sqrt{2}}
\]

\[
\therefore \ c_0 = c_1 = \frac{1}{\sqrt{2}} .
\]

Dari perhitungan didapat solusi yang konsisten, yaitu nilai
$a_0 = 1$, $a_1 = 0$, dan $b_0 = b_1 = c_0 = c_1 = \frac{1}{\sqrt{2}}$, sehingga keadaan $|\chi_1\rangle$ dapat dinyatakan menjadi

\begin{equation}
|\chi_1\rangle
=
|0\rangle_A \otimes
\left(
\frac{1}{\sqrt{2}}
(|0\rangle + |1\rangle)
\right)_B
\otimes
\left(
\frac{1}{\sqrt{2}}
(|0\rangle + |1\rangle)
\right)_C .
\end{equation}

\subsubsection{Keadaan Terbelit Sepenuhnya}

Semisal ditinjau keadaan tiga qubit yang lain, semisal keadaan $|\chi_2\rangle$ sebagai berikut

\begin{equation}
|\chi_2\rangle =
\frac{1}{\sqrt{3}}
\left(
|001\rangle + |010\rangle + |100\rangle
\right) .
\end{equation}

Dengan cara yang sama seperti diatas, kita bisa mencari masing-masing koefisien keadaan qubit tunggalnya

\begin{equation}
\frac{a_0 b_0 c_0}{a_1 b_0 c_0}
=
\frac{0}{1/\sqrt{3}}
\end{equation}

\begin{equation}
\frac{a_0 b_0 c_0}{a_0 b_1 c_0}
=
\frac{0}{1/\sqrt{3}}
\end{equation}

\begin{equation}
\frac{a_0 b_0 c_0}{a_0 b_0 c_1}
=
\frac{0}{1/\sqrt{3}}
\end{equation}

\[
\therefore \ a_0 = 0 ; \ a_1 = 1
\qquad
\therefore \ b_0 = 0 ; \ b_1 = 1
\qquad
\therefore \ c_0 = 0 ; \ c_1 = 1
\]

\begin{equation}
\frac{a_1 b_0 c_1}{a_0 b_0 c_1}
=
\frac{0}{1/\sqrt{3}}
\end{equation}

\begin{equation}
\frac{a_0 b_1 c_1}{a_0 b_0 c_1}
=
\frac{0}{1/\sqrt{3}}
\end{equation}

\begin{equation}
\frac{a_0 b_1 c_1}{a_0 b_1 c_0}
=
\frac{0}{1/\sqrt{3}}
\end{equation}

\[
\therefore \ a_1 = 0 ; \ a_0 = 1
\qquad
\therefore \ b_1 = 0 ; \ b_0 = 1
\qquad
\therefore \ c_1 = 0 ; \ c_0 = 1 .
\]

Terlihat bahwa nilai $a_0$ dan $a_1$ tidak konsisten, begitu juga dengan nilai $b_0$, $b_1$, $c_0$, dan $c_1$. Karena tidak ada solusi yang pasti untuk nilai koefisien-koefisien qubit tunggal, maka keadaan $|\chi_2\rangle$ ini dinamakan keadaan terbelit sepenuhnya (\textit{completely entangled state}).

\subsubsection{Keadaan Gabungan}

Ditinjau contoh keadaan lainnya $|\chi_3\rangle$ yang cukup berbeda dari contoh-contoh keadaan sebelumnya

\begin{equation}
|\chi\rangle_3 =
\frac{1}{\sqrt{2}}
\left(
|000\rangle + |011\rangle
\right) ,
\end{equation}

dengan masing-masing nilai koefisien keadaan qubit tunggal nya adalah sebagai berikut

\begin{equation}
\frac{a_1 b_0 c_0}{a_0 b_0 c_0}
=
\frac{0}{1/\sqrt{2}}
\end{equation}

\begin{equation}
\frac{a_0 b_1 c_0}{a_0 b_0 c_0}
=
\frac{0}{1/\sqrt{2}}
\end{equation}

\begin{equation}
\frac{a_0 b_0 c_1}{a_0 b_0 c_0}
=
\frac{0}{1/\sqrt{2}}
\end{equation}

\[
\therefore \ a_1 = 0 ; \ a_0 = 1
\qquad
\therefore \ b_1 = 0 ; \ b_0 = 1
\qquad
\therefore \ c_1 = 0 ; \ c_0 = 1
\]

\begin{equation}
\frac{a_1 b_1 c_1}{a_0 b_1 c_1}
=
\frac{0}{1/\sqrt{2}}
\end{equation}

\begin{equation}
\frac{a_0 b_0 c_1}{a_0 b_1 c_1}
=
\frac{0}{1/\sqrt{2}}
\end{equation}

\begin{equation}
\frac{a_0 b_0 c_1}{a_0 b_1 c_1}
=
\frac{0}{1/\sqrt{2}}
\end{equation}

\[
\therefore \ a_1 = 0 ; \ a_0 = 1
\qquad
\therefore \ b_0 = 0 ; \ b_1 = 1
\qquad
\therefore \ c_0 = 0 ; \ c_1 = 1 .
\]

Terlihat bahwa nilai $a_0$ dan $a_1$ mempunyai solusi yang pasti, namun nilai $b_0$ dan $b_1$ tidak mempunyai solusi yang pasti, begitu juga dengan nilai $c_0$ dan $c_1$. Hal ini tidak sesuai dengan syarat keadaan terpisah sepenuhnya karena solusi untuk nilai koefisien keadaan sub-sistem $B$ ($b_0$ dan $b_1$) dan nilai koefisien sub-sistem $C$ ($c_0$ dan $c_1$) tidak dapat dicari secara pasti. Akan tetapi, keadaan ini juga tidak sesuai dengan syarat keadaan terbelit sepenuhnya karena solusi untuk nilai koefisien keadaan sub-sistem $A$ ($a_0$ dan $a_1$) dapat dicari secara pasti. Maka dari itu, keadaan seperti ini dinamakan
keadaan gabungan (\textit{compound state}) (Purwanto, Sukamto, \& Yuwana, 2018).
Keadaan $|\chi_3\rangle$ dinamakan keadaan gabungan $A$-$BC$
karena keadaan sistem $ABC$ tersebut tersusun atas sub-sistem $A$ dan sub-sistem $BC$ yang saling terpisah, namun keadaan sub-sistem $BC$ merupakan keadaan yang terbelit. Keadaan gabungan $A$-$BC$ ini bisa dituliskan kembali menjadi

\begin{equation}
|\chi\rangle_3
=
|\psi\rangle_{A-BC}
=
|0\rangle_A \otimes
\left(
\frac{1}{\sqrt{2}}
(|00\rangle + |11\rangle)
\right)_{BC} .
\end{equation}

Contoh keadaan gabungan yang lain adalah keadaan gabungan $AB$-$C$ yang bisa dituliskan sebagai berikut

\begin{equation}
|\psi\rangle_{AB-C}
=
\frac{1}{2}
\left(
|001\rangle + |011\rangle + |101\rangle - |111\rangle
\right) .
\end{equation}

Jika dilakukan perhitungan untuk mencari solusi nilai koefisien keadaan qubit tunggal nya dengan cara yang sama seperti sebelumnya, maka bisa didapatkan bahwa nilai koefisien
$a_0$, $a_1$, $b_0$, dan $b_1$ tidak mempunyai solusi yang pasti, namun nilai koefisien $c_0$ dan $c_1$ bisa dicari dengan pasti, yaitu $c_0 = 0$ dan $c_1 = 1$, sehingga keadaan $|\psi\rangle_{AB-C}$ bisa dituliskan ulang menjadi

\begin{equation}
|\psi\rangle_{AB-C}
=
\left(
\frac{1}{2}
(|00\rangle + |01\rangle + |10\rangle - |11\rangle)
\right)_{AB}
\otimes
|1\rangle_C .
\end{equation}

Selain keadaan gabungan $A$-$BC$ (sub-sistem $A$ dan sub-sistem $BC$) dan $AB$-$C$ (sub-sistem $AB$ dan sub-sistem $C$), terdapat satu keadaan gabungan lainnya yaitu keadaan gabungan $AC$-$B$ (sub-sistem $AC$ dan sub-sistem $B$), dengan contoh
keadaannya bisa dituliskan sebagai

\begin{equation}
|\psi\rangle_{AC-B}
=
\frac{1}{2}
\left(
|000\rangle + |010\rangle + |101\rangle + |111\rangle
\right) .
\end{equation}

Kita bisa membuktikan bahwa sub-sistem $A$ dan $C$ terbelit, namun sub-sistem $B$ terpisah, menggunakan cara yang sama seperti kita mencari koefisien qubit tunggal, kita bisa dapatkan bahwa nilai $a_0$ dan $a_1$ tidak konsisten sehingga tidak mempunyai solusi, begitu juga dengan nilai $c_0$ dan $c_1$, namun untuk nilai $b_0$ dan $b_1$ konsisten sehingga mempunyai solusi

\[
a_0 = 0 ; \ a_1 = 1
\qquad
\text{dan}
\qquad
a_1 = 0 ; \ a_0 = 1
\]

\[
b_0 = b_1 = \frac{1}{\sqrt{2}}
\]

\[
c_0 = 0 ; \ c_1 = 1
\qquad
\text{dan}
\qquad
c_1 = 0 ; \ c_0 = 1 .
\]

Namun, kita tidak bisa menuliskan keadaan $|\psi\rangle_{AC-B}$ seperti berikut

\begin{equation}
|\psi\rangle_{AC-B}
\neq
\left(
\frac{1}{\sqrt{2}}
(|00\rangle + |11\rangle)
\right)_{AC}
\otimes
\left(
\frac{1}{\sqrt{2}}
(|0\rangle + |1\rangle)
\right)_B ,
\end{equation}

karena justru akan menjadi keadaan gabungan $AB$-$C$. Kita juga tidak bisa menuliskan keadaan menjadi

\begin{equation}
|\psi\rangle_{AC-B}
\neq
\left(
\frac{1}{\sqrt{2}}
(|0\rangle + |1\rangle)
\right)_B
\otimes
\left(
\frac{1}{\sqrt{2}}
(|00\rangle + |11\rangle)
\right)_{AC} ,
\end{equation}

karena justru akan menjadi keadaan gabungan $A$-$BC$.

\subsection{Rank Matriks Densitas Tereduksi}

\noindent
\begin{table}[h]
\centering
\caption{Matriks Densitas Dua Qubit}
\begin{tabular}{|c|c|c|c|c|c|c|}
    \hline
    Keadaan Dua Qubit & $r(\rho_A)$ & $r(\rho_B)$ & $r(\rho_C)$ & $r(\rho_{AB})$ & $r(\rho_{AC})$ & $r(\rho_{BC})$ \\
    \hline
    Terpisah $\rho^{(sep)}$ & 1 & 1 & -- & -- & -- & -- \\
    \hline
    Terbelit $\rho^{(ent)}$ & 2 & 2 & -- & -- & -- & -- \\
    \hline
\end{tabular}
\end{table}

\vspace{0.5cm}

\noindent
\begin{table}[h]
\centering
\caption{Matriks Densitas Tiga Qubit}
\begin{tabular}{|c|c|c|c|c|c|c|}
    \hline
    \textbf{Keadaan Tiga Qubit} & $r(\rho_A)$ & $r(\rho_B)$ & $r(\rho_C)$ & $r(\rho_{AB})$ & $r(\rho_{AC})$ & $r(\rho_{BC})$ \\
    \hline
    Terpisah sepenuhnya $\rho^{(sep)}$ & 1 & 1 & 1 & 1 & 1 & 1 \\
    \hline
    Terbelit sepenuhnya $\rho^{(ent)}$ & 2 & 2 & 2 & 2 & 2 & 2 \\
    \hline
    Gabungan A-BC & 1 & 2 & 2 & 2 & 2 & 2 \\
    \hline
    Gabungan AB-C & 2 & 2 & 1 & 2 & 2 & 2 \\
    \hline
    Gabungan AC-B & 2 & 1 & 2 & 2 & 2 & 2 \\
    \hline
\end{tabular}
\end{table}

\par
Berdasarkan perumusan bentuk umum matriks densitas tereduksi beserta nilai rank-nya yang dirangkum dalam Tabel 2.1 dan 2.2, dapat diamati bahwa untuk keadaan dua qubit yang saling terpisah, seluruh matriks densitas tereduksi $\rho_\kappa$ ($\kappa$ menyatakan masing-masing sub-sistem) memiliki rank sebesar 1. Sebaliknya, pada keadaan terbelit, setiap $\rho_\kappa$ memiliki rank sebesar 2 \cite{nielsen2010quantum,horodecki2009quantum}. Fenomena yang sepadan juga dijumpai pada sistem tiga qubit, di mana keadaan yang sepenuhnya terpisah ditandai oleh rank $\rho_\kappa = 1$ untuk seluruh sub-sistem, sedangkan pada keadaan terbelit sepenuhnya, rank masing-masing matriks densitas tereduksi bernilai 2 \cite{horodecki2009quantum}. Untuk keadaan gabungan, misalnya konfigurasi A–BC, diperoleh bahwa rank $\rho_A$ sama dengan rank $\rho_{BC}$, yaitu bernilai 1, sementara matriks densitas tereduksi sub-sistem lainnya memiliki rank sebesar 2 \cite{dur2000three}. Dengan demikian, suatu keadaan kuantum dapat dikatakan terpisah jika dan hanya jika matriks densitas tereduksinya mempunyai rank sama dengan 1 \cite{nielsen2010quantum,bengtsson2017geometry}. Sebagai ilustrasi, pada keadaan dua qubit yang terpisah, keadaan masing-masing sub-sistem A dan B dapat ditentukan secara pasti. Oleh karena itu, kedua sub-sistem tersebut dapat diperlakukan sebagai sistem satu qubit yang berdiri sendiri dan tidak saling berinteraksi, sehingga matriks densitas masing-masing sub-sistem dapat dikonstruksi secara independen sebagaimana matriks densitas keadaan awal \cite{nielsen2010quantum}.
\par
Pada sistem tiga qubit, khususnya untuk keadaan gabungan, struktur keadaan dapat dipahami sebagai pemisahan antara satu sub-sistem satu qubit dan satu sub-sistem dua qubit yang berada dalam keadaan terbelit \cite{dur2000three}. Misalnya pada keadaan AB-C, diperoleh bahwa rank $\rho_A$ dan rank $\rho_C$ sama-sama bernilai 1, yang menunjukkan bahwa sub-sistem AB terpisah dari sub-sistem C. Akibatnya, kedua sub-sistem tersebut dapat direpresentasikan secara terpisah menggunakan matriks densitas asal masing-masing \cite{dur2000three}. Sebagai contoh lain, nilai rank $\rho_{AC}$ adalah 2. Meskipun keadaan sub-sistem C dapat diketahui secara pasti karena berada dalam keadaan terpisah, keadaan sub-sistem A tidak dapat ditentukan secara pasti karena masih merupakan bagian dari sub-sistem terbelit AB, sehingga matriks densitas tereduksinya memiliki rank sebesar 2, atau secara lebih khusus tidak sama dengan 1 \cite{horodecki2009quantum,bengtsson2017geometry}.

\subsection{Entropi Von Neumann}

\begin{table}[h]
\centering
\caption{Entropi Von Neumann Dua dan Tiga Qubit}
\begin{tabular}{|c|c|c|c|c|c|c|}
\hline
 & $r(\rho_A)$ & $r(\rho_B)$ & $r(\rho_C)$ & $r(\rho_{AB})$ & $r(\rho_{AC})$ & $r(\rho_{BC})$ \\
\hline
\textbf{Keadaan Dua Qubit} & & & & & & \\
\hline
Terpisah $S(\rho^{(sep)})$ & 0 & 0 & -- & -- & -- & -- \\
\hline
Terbelit $S(\rho^{(ent)})$ & $>$ 0 & $>$ 0 & -- & -- & -- & -- \\
\hline
\textbf{Keadaan Tiga Qubit} & & & & & & \\
\hline
Terpisah sepenuhnya $S(\rho^{(sep)})$ & 0 & 0 & 0 & 0 & 0 & 0 \\
\hline
Terbelit sepenuhnya $S(\rho^{(ent)})$ & $>$ 0 & $>$ 0 & $>$ 0 & $>$ 0 & $>$ 0 & $>$ 0 \\
\hline
Gabungan A-BC & 0 & $>$ 0 & $>$ 0 & $>$ 0 & $>$ 0 & 0 \\
\hline
Gabungan AB-C & $>$ 0 & $>$ 0 & 0 & 0 & $>$ 0 & $>$ 0 \\
\hline
Gabungan AC-B & $>$ 0 & 0 & $>$ 0 & $>$ 0 & 0 & $>$ 0 \\
\hline
\end{tabular}
\end{table}

\par
Berdasarkan perumusan bentuk umum matriks densitas tereduksi beserta nilai entropi Von Neumann yang dirangkum dalam Tabel 2.3, diperoleh bahwa pada sistem dua qubit yang berada dalam keadaan terpisah, entropi masing-masing sub-sistem $S^{(sep)}$ (dengan A menyatakan setiap sub-sistem) bernilai nol. Sebaliknya, pada keadaan terbelit, entropi sub-sistem $S^{(ent)}$ bernilai lebih besar dari nol \cite{nielsen2010quantum,wehrl1978general}. Pola yang serupa juga ditemukan pada sistem tiga qubit, di mana keadaan yang sepenuhnya terpisah ditandai oleh entropi nol pada seluruh sub-sistem, sedangkan pada keadaan terbelit sepenuhnya setiap sub-sistem memiliki entropi bernilai positif \cite{horodecki2009quantum}.
\par
Pada keadaan gabungan, misalnya konfigurasi AB-C, terlihat bahwa entropi $S_{C}^{(sep)}$ memiliki nilai yang sama dengan entropi $S_{AB}^{(sep)}$, yaitu nol, sementara entropi sub-sistem lainnya bernilai lebih besar dari nol \cite{dur2000three}. Hal ini menunjukkan bahwa suatu keadaan kuantum dapat dikategorikan sebagai keadaan terpisah jika dan hanya jika entropi sub-sistemnya bernilai nol \cite{nielsen2010quantum,bengtsson2017geometry}. Sebagai contoh, pada keadaan dua qubit yang terpisah, keadaan masing-masing sub-sistem A dan B dapat ditentukan secara pasti. Oleh karena itu, masing-masing sub-sistem tersebut dapat dipandang sebagai sistem qubit tunggal yang berada dalam keadaan murni dan dapat direpresentasikan oleh vektor keadaan yang terdefinisi dengan baik, sehingga tidak mengandung ketidakpastian tambahan dan menghasilkan nilai entropi nol \cite{wehrl1978general}.
\par
Pada sistem tiga qubit, khususnya untuk keadaan gabungan, struktur keadaan dapat dipahami sebagai pemisahan antara satu sub-sistem qubit tunggal dan satu sub-sistem dua qubit yang berada dalam keadaan terbelit \cite{dur2000three}. Misalnya pada keadaan AB$|$C, diperoleh bahwa entropi $S_{C}^{(sep)}$ dan entropi $S_{AB}^{(sep)}$ sama-sama bernilai nol. Hal ini menandakan bahwa sub-sistem AB dan sub-sistem C masing-masing berada dalam keadaan murni yang tidak saling berkorelasi, sehingga keduanya dapat direpresentasikan secara terpisah melalui matriks densitas asal dan diklasifikasikan sebagai keadaan terpisah \cite{horodecki2009quantum}.
\par
Sebagai ilustrasi lain, dapat ditinjau entropi $S_{AC}^{(cp2)}$ yang bernilai lebih besar dari nol. Meskipun keadaan sub-sistem C dapat diketahui secara pasti karena berada dalam keadaan terpisah, keadaan sub-sistem A tidak dapat ditentukan secara pasti karena masih merupakan bagian dari sub-sistem terbelit AB. Akibatnya, keadaan sub-sistem A memiliki ketidakpastian dalam representasi vektor keadaannya di ruang Hilbert. Oleh karena itu, ketika sub-sistem AC dipertimbangkan secara keseluruhan, keadaan ini tidak dapat direpresentasikan sebagai vektor keadaan murni yang tunggal, sehingga menghasilkan tambahan informasi mengenai kemungkinan keadaan yang ada dan menyebabkan entropi $S_{AC}^{(cp2)}$ bernilai lebih besar dari nol \cite{nielsen2010quantum,horodecki2009quantum}.
